\section{Dynamical CBED: Many-Beam Coupling, Absorption, HOLZ Lines, and Point-Group
Determination}

The companion note \texttt{docs/tex/algorithms/convergent\_beam\_electron\_diffraction.tex}
derives the geometry of a convergent-beam pattern and the two-beam rocking curve inside one
disc. That treatment is deliberately incomplete in four connected ways, and this note closes
them: the discs of one pattern are computed independently and so are not mutually consistent;
nothing is absorbed, so the fringes never decay; the higher-order Laue zones appear only as
ring radii and not as the sharp lines used for metrology; and the symmetry analysis that
yields the crystal point group --- including the presence or absence of a centre of
symmetry --- is absent altogether.

Those four are one capability rather than four. Many-beam coupling makes the pattern a
single object; absorption makes it a physically realizable one; admitting higher-order Laue
zone beams into the same beam set both produces the HOLZ lines and destroys the projection
symmetry that would otherwise make every pattern look centrosymmetric; and the symmetry that
survives that destruction \emph{is} the diffraction group. The implementation is
\texttt{pytex.diffraction.dynamical}, \texttt{pytex.diffraction.holz} and
\texttt{pytex.diffraction.diffraction\_groups}.

\subsection{The Coupled Beam Equations}

Write the wavefield inside the crystal as a sum over reciprocal-lattice beams,
\(\Psi(\mathbf{r}) = \sum_{g}\psi_{g}(z)\exp[i(\mathbf{k}+\mathbf{g})\cdot\mathbf{r}]\).
Substituting into the Schr\"odinger equation for a fast electron in the crystal potential and
neglecting backscattering (the column approximation) gives, for a parallel-sided slab,

\begin{equation}
  \frac{\mathrm{d}\psi}{\mathrm{d}z} = i\pi\,\mathbf{A}\,\psi,
  \qquad
  A_{gg} = 2 s_{g},
  \qquad
  A_{gh} = \nu_{g-h} \quad (g \neq h),
  \label{eq:dyn-system}
\end{equation}

with the excitation errors \(s_{g}(\boldsymbol{\theta})\) of Eq.~(1) of the companion note
and the Fourier coefficients of the scaled potential

\begin{equation}
  \nu_{g} = \frac{\lambda F_{g}}{\pi V_{c}\cos\theta_{g}},
  \qquad |\nu_{g}| = \frac{1}{\xi_{g}} .
  \label{eq:dyn-nu}
\end{equation}

The scaling in \eqref{eq:dyn-nu} is chosen so that the modulus of the coupling coefficient is
exactly the reciprocal of the two-beam extinction distance already validated against Williams
and Carter's Table 23.1. The absolute scale of the many-beam calculation is therefore not a
second, independent claim: it is the same claim, and it is pinned by the two-beam limit below.

Because \(\mathbf{A}\) does not depend on \(z\), \eqref{eq:dyn-system} integrates in closed
form,

\begin{equation}
  \psi(t) = \exp(i\pi\mathbf{A}t)\,\mathbf{e}_{0},
  \label{eq:dyn-propagator}
\end{equation}

the incident boundary condition being unit amplitude in the transmitted beam. Writing
\(\mathbf{A} = \mathbf{C}\,\mathrm{diag}(\gamma_{j})\,\mathbf{C}^{-1}\),

\begin{equation}
  \psi_{g}(t) = \sum_{j} C_{gj}\,\alpha_{j}\,e^{i\pi\gamma_{j}t},
  \qquad
  \boldsymbol{\alpha} = \mathbf{C}^{-1}\mathbf{e}_{0},
  \qquad
  I_{g} = |\psi_{g}|^{2}.
  \label{eq:dyn-bloch}
\end{equation}

The eigenvectors are the \emph{Bloch waves}: wavefields that propagate unchanged in shape,
each attenuated at its own rate \(\mathrm{Im}\,\gamma_{j}\). The real parts of the
\(\gamma_{j}\) trace the dispersion surface.

\paragraph{Why the excitation amplitudes are solved for and not projected.} \(\mathbf{A}\) is
Hermitian without absorption and neither Hermitian nor normal with it, so \(\mathbf{C}\) is
not orthogonal in general. Obtaining \(\boldsymbol{\alpha}\) by projection, \(\alpha_{j} =
C_{0j}^{*}\), is the classic silent error in a Bloch-wave implementation: it produces
rocking curves of the right shape with the wrong contrast, and it breaks intensity
conservation in a way that is only visible if one checks for it.

\subsection{Three Exact Properties}

\paragraph{The two-beam limit.} For a single reflection, \(\mathbf{A} = s\mathbf{I} +
\mathbf{B}\), where \(\mathbf{B}\) is the traceless matrix with diagonal \((-s, +s)\) and
both off-diagonal entries equal to \(\nu = 1/\xi_{g}\). It has eigenvalues
\(\pm s_{\mathrm{eff}}\), \(s_{\mathrm{eff}} =
\sqrt{s^{2}+\xi_{g}^{-2}}\), and the exponential of a traceless \(2\times 2\) matrix is
\(\cos(\pi s_{\mathrm{eff}}t)\mathbf{I} + i\sin(\pi s_{\mathrm{eff}}t)\mathbf{B}/
s_{\mathrm{eff}}\), so

\begin{equation}
  I_{g} = \frac{\sin^{2}(\pi t s_{\mathrm{eff}})}{(\xi_{g}s_{\mathrm{eff}})^{2}},
\end{equation}

which is the Howie--Whelan expression the companion note derives. This is asserted to
machine precision, and it pins the diagonal convention, the off-diagonal scale and the
factor \(i\pi\) simultaneously.

\paragraph{Unitarity.} Without absorption \(\mathbf{A}\) is Hermitian, so
\(\exp(i\pi\mathbf{A}t)\) is unitary and \(\sum_{g}I_{g} = 1\) at every thickness and every
incident direction. This is the only exact global check available on a many-beam
calculation, and it is lost the moment absorption is switched on --- which is a reason to
run it once with absorption off.

\paragraph{Normal absorption is a scalar.} The mean absorptive coefficient
\(i/\xi'_{0}\) sits on \emph{every} diagonal element, so it commutes out of the exponential:

\begin{equation}
  \exp\left[i\pi\left(\mathbf{A} + \tfrac{i}{\xi'_{0}}\mathbf{I}\right)t\right]
    = e^{-\pi t/\xi'_{0}}\exp(i\pi\mathbf{A}t),
  \qquad
  I_{g} \to e^{-2\pi t/\xi'_{0}}\,I_{g}.
  \label{eq:dyn-normal}
\end{equation}

It can therefore change no relative intensity, no fringe position and no symmetry. Every
qualitative effect of absorption comes from the \emph{off-diagonal} absorptive terms.

\subsection{Absorption as an Imaginary Optical Potential}

Electrons leave the coherent elastic wavefield by processes the calculation does not follow:
thermal diffuse scattering above all, then plasmon and core losses. The standard
representation adds an imaginary part to each Fourier coefficient,

\begin{equation}
  \frac{1}{\xi_{g}} \;\longrightarrow\; \frac{1}{\xi_{g}} + \frac{i}{\xi'_{g}},
\end{equation}

so that \eqref{eq:dyn-system} becomes \(A_{gg} = 2s_{g} + i/\xi'_{0}\) and \(A_{gh} =
\nu_{g-h}(1 + i r)\) with \(r = \xi_{g}/\xi'_{g}\).

\paragraph{What is derived and what is assumed.} The \emph{structure} is not an
approximation: an imaginary optical potential is the correct representation of loss from a
coherent wavefield, and the observable consequences below follow from the eigenvector
structure rather than being applied afterwards. The \emph{magnitudes} are phenomenological.
\texttt{pytex} carries them as the two ratios \(\xi_{0}/\xi'_{0}\) and \(\xi_{g}/\xi'_{g}\),
whose customary working value for a metal near 100--200~kV is about
\(0.1\) (Hirsch \emph{et al.}, \emph{Electron Microscopy of Thin Crystals}, Ch.~12). A
first-principles absorptive form factor --- the Einstein-model thermal-diffuse integral of
Hall and Hirsch, as parametrized by Bird and King --- is not implemented, and the
documentation says so rather than implying a computed constant.

\paragraph{Positivity.} The absorptive matrix has eigenvalues \(\xi_{0}'^{-1} \pm
\xi_{g}'^{-1}\) in the two-beam case. Since \(\xi_{0} < \xi_{g}\) always, requiring the
reflection ratio not to exceed the mean ratio is sufficient for both to be positive. A
larger reflection ratio would produce a Bloch wave that \emph{gains} intensity with depth,
which no absorption process can do; the constructor refuses it.

\paragraph{Anomalous absorption.} The two Bloch waves of a two-beam calculation have their
intensity maxima on and between the atomic planes respectively. The first is absorbed
strongly, the second weakly, and which is preferentially excited depends on the sign of
\(s\). The consequence is the Hashimoto--Howie--Whelan result: with absorption, the
bright-field rocking curve becomes asymmetric about \(s = 0\) while the dark-field one
remains exactly symmetric. \texttt{pytex} asserts both halves of that statement numerically;
they are a stronger test of the absorption implementation than any single number, because
they are properties no incorrect implementation reproduces by accident.

\subsection{Friedel's Law as a Propagator Theorem}

For a real potential \(\nu_{-g} = \nu_{g}^{*}\), so \(\mathbf{A}\) is always Hermitian. It is
\emph{symmetric} only when every included \(\nu_{g}\) is real, which happens exactly when the
structure sampled by the beam set has a centre of symmetry with the origin on it.

Relabelling the beams by \(g \mapsto -g\) and the incident tilt by \(\boldsymbol{\theta}
\mapsto -\boldsymbol{\theta}\) maps the diagonal \(2s_{g}(\boldsymbol{\theta}) \mapsto
2s_{-g}(-\boldsymbol{\theta})\) and the off-diagonal \(\nu_{g-h} \mapsto \nu_{h-g} =
\nu_{g-h}^{*}\). In the zeroth Laue zone \(s_{-g}(-\boldsymbol{\theta}) =
s_{g}(\boldsymbol{\theta})\) exactly, so the relabelled matrix is
\(\mathbf{A}^{\mathsf{T}}\), the propagator becomes \(M^{\mathsf{T}}\), and

\begin{equation}
  I_{g}(\boldsymbol{\theta}) = I_{-g}(-\boldsymbol{\theta})
  \iff |M_{g0}| = |M_{0g}|
  \iff \mathbf{M} \text{ symmetric}
  \iff \text{the sampled structure is centrosymmetric.}
  \label{eq:dyn-friedel}
\end{equation}

Friedel's law is thus recovered as a theorem about the propagator rather than as a kinematic
accident, and \eqref{eq:dyn-friedel} is the mechanism behind the point-group determination.

\paragraph{The word ``sampled'' carries the whole difficulty.} A beam set confined to the
zeroth Laue zone samples the potential \emph{projected} along the beam, and that projection
is frequently centrosymmetric when the crystal is not. Zincblende down \([111]\) is the
standard example: every ZOLZ coefficient is real, so a projection calculation reports
Friedel's law to machine precision and cannot see the polarity of the crystal. Admitting the
first-order Laue zone gives the coefficients phases, breaks the symmetry of the propagator,
and separates the two discs of a \(\pm\mathbf{g}\) pair by tens of percent. \texttt{pytex}
pins this with a controlled pair: zincblende GaAs and a rocksalt structure built on the same
lattice from the same two species, differing only by the position of the second sublattice
and therefore only by the centre of symmetry. The Friedel violation separates them by more
than three orders of magnitude.

\emph{Higher-order Laue zone interaction is not a refinement here; it is the entire
mechanism.} A symmetry conclusion drawn from a zeroth-Laue-zone calculation is worthless, and
\texttt{BeamSet.holz\_mask} exists so that this can be checked rather than assumed.

\subsection{Beam Selection}

A many-beam calculation is defined only once the beams are named. \texttt{pytex} selects a
reflection when it comes within a stated window of the Bragg condition \emph{for some
incident direction in the illumination cone}. Because \(s_{g}\) is affine in the tilt,

\begin{equation}
  \min_{|\boldsymbol{\theta}| \le \alpha} |s_{g}(\boldsymbol{\theta})|
    = \max\bigl(0,\; |s_{g}(\mathbf{0})| - \alpha|\mathbf{g}_{\perp}|\bigr),
  \label{eq:dyn-selection}
\end{equation}

which is evaluated in closed form. The distinction matters: a HOLZ reflection is far from
Bragg at the centre of the pattern and exactly at Bragg somewhere inside the bright-field
disc --- that locus is a HOLZ line --- so selecting on the zero-tilt excitation error alone
would discard every HOLZ reflection and, with it, the symmetry-breaking mechanism of the
previous subsection.

\paragraph{Cost.} Every selected beam is solved exactly; there is no Bethe perturbation of
weak beams, so the cost is \(O(m n^{3})\) for \(m\) incident directions and \(n\) beams. The
economy that works is a tighter window in \eqref{eq:dyn-selection}, which removes beams that
were barely coupled, rather than coarser tilt sampling, which degrades the fringe positions
the calculation exists to predict.
